% !PRINTED 549
\weberpage{549}
% running head: \S.\,151. Ideale. 549
%<<CONT
assign. The same ideal $\mathfrak{a}$ is then also assigned to all the
functionals associated with $\varphi$. This assignment has the
properties 1., 4., 5.

If $\varphi$ and $\varphi_1$ are two non-associated functionals, then certainly
one of them, say $\varphi$, is not divisible by the other $\varphi_1$,
and consequently (by \S.\,143, 3.) there are integers that are
divisible by $\varphi$ but not by $\varphi_1$. Consequently, \gesperrt{different}
ideals $\mathfrak{a}$, $\mathfrak{a}_1$ are always assigned to non-associated functionals.

But it must now also be shown that in this way \gesperrt{all}
ideals of the field $\Omega$ can be obtained; in other words, that
every ideal $\mathfrak{a}$ different from $\mathfrak{o}$ consists of the totality
of the numbers divisible by a certain functional factor.

We therefore now start from any ideal $\mathfrak{a}$ and
choose from it an arbitrary finite set of numbers,
$\alpha_1, \alpha_2, \ldots, \alpha_r$, whose greatest common divisor may be
$\delta_r$. This functional $\delta_r$ has a finite number of prime
factors.

Now if there is a number $\alpha_{r+1}$ in $\mathfrak{a}$ that is not divisible
by $\delta_r$, then the greatest common divisor $\delta_{r+1}$ of
$\delta_r$ and $\alpha_{r+1}$ has fewer prime factors than $\delta_r$. If we
continue with this mode of reasoning, we come to the result that
a finite number of numbers $\alpha_1, \alpha_2, \ldots, \alpha_m$ can be chosen
from $\mathfrak{a}$ in such a way that the greatest common divisor $\delta$ of these
numbers \gesperrt{divides all the numbers of $\mathfrak{a}$.}

On the other hand, every number in $\mathfrak{o}$ divisible by $\delta$ belongs to $\mathfrak{a}$.
For by \S.\,149, 4. every number $\alpha$ divisible by $\delta$ can be
put in the form

\[
\alpha = \alpha_1 \xi_1 + \alpha_2 \xi_2 + \cdots + \alpha_m \xi_m,
\]

in which $\xi_1, \xi_2, \ldots, \xi_m$ are numbers in $\mathfrak{o}$, and consequently,
by I. and II., $\alpha$ belongs to the ideal $\mathfrak{a}$.

From this it results that, when $\delta$ is a unit, the
ideal $\mathfrak{a}$ is identical with $\mathfrak{o}$. Every ideal $\mathfrak{a}$ is thus
characterised by the fact that all its numbers have a certain greatest
common divisor.

It remains still to be shown that, when the functionals $\varphi, \psi$
correspond to the two ideals $\mathfrak{a}, \mathfrak{b}$, the product $\varphi\psi$
corresponds to the ideal $\mathfrak{a}\mathfrak{b}$.
